\begin{exercisechunk}
\item \textbf{Characterize causally orderable self-maps.}\suppmark
\label[exercise]{ex:l14-causal-order}
\exerciseprofile{2}{2}{2}{\exproof\quad\exconstruction\quad\exinterpretation}{%
Characterize exactly when every pointer update can retrieve from an earlier
record in one causal ordering.}
The pointer-doubling construction requires an ordering in which every map
record points backward or to itself. This exercise identifies exactly which
finite self-maps admit such an ordering.

Let \(p:\cD\to\cD\) be a self-map on a finite set. Its functional directed
graph has one edge \(u\to p(u)\) from every vertex. A \emph{causal ordering}
is a bijection \(\operatorname{pos}:\cD\to[|\cD|]\) satisfying
\[
 \operatorname{pos}(p(u))\le\operatorname{pos}(u)
 \qquad(u\in\cD).
\]
\begin{enumerate}[(a),beginpenalty=10000]
\item Prove that if the functional graph contains a directed cycle of length
greater than one, then no causal ordering exists.
\item Suppose every directed cycle is a self-loop. Prove that every orbit
eventually reaches a unique fixed point in its connected component. Define the
depth of a vertex to be the number of map applications needed to reach that
fixed point. Construct a causal ordering and prove that it satisfies the
required inequality.
\item Conclude that a causal ordering exists if and only if every directed
cycle is a self-loop.
\item Prove that a causal ordering also satisfies
\[
 \operatorname{pos}(p^k(u))\le\operatorname{pos}(u)
 \qquad(u\in\cD,\ k\ge0).
\]
Explain why this inequality makes every lookup in causal pointer-doubling
preprocessing available at or before the requesting record.
\end{enumerate}
\end{exercisechunk}
